\scanpage{43}
\mspage{- 24 -}
% archive sheet number: 42

\[
\left\{\begin{aligned}
m-1&=n+\handwritten{\sum_{j=1}^p\sum_{k=1}^{r_j}}s_{kj}\\
\handwritten{\mu_\ell}&=\handwritten{\unsure{\lambda|\Lambda}_\ell}+\handwritten{\sum_{j=1}^p\sum_{k=1}^{r_j}\nu}_{kj}\qquad(1\handwritten{\leq\ell\leq}q)
\end{aligned}\right.\tag*{(16)}
\]
On tire de ces conditions que
\[
m+\handwritten{\sum_{\ell=1}^q\mu_\ell}=1+n+\handwritten{\sum_{\ell=1}^q\lambda_\ell}+\handwritten{\sum_{j=1}^p\sum_{k=1}^{r_j}}(s_{kj}+\handwritten{\sum_{\ell=1}^q\nu}_{kj\handwritten{\ell}})
\]
et par suite, pour tout système d'exposants $(s_{kj},\handwritten{\nu}_{kj1},..,\handwritten{\nu}_{kjq})$ qui figure aux seconds membres de (16), on a $s_{kj}+\handwritten{\sum_{\ell=1}^q\nu}_{kj\handwritten{\ell}}\handwritten{<}m+\handwritten{\sum_{\ell=1}^q\mu_\ell}$ ; en d'autres termes, et comme K est de caractéristique 0, les relations (15) déterminent les coefficients des termes de degré \uline{total} t dans $u_i$ en fonction de ceux des termes de degré total $\handwritten{<}t$ dans les $u_j$ $(1\handwritten{\leq}j\handwritten{\leq}p)$ ; par récurrence sur t, on voit donc que les équations (15) déterminent les coefficients $h_i(m;\handwritten{\mu}_1,..,\handwritten{\mu}_q)$ de fa\handwritten{-}çon \uline{unique}, ce qui montre déjà que si le problème est possible, il n'a qu'une seule solution. En outre, si les $h_i(m;\handwritten{\mu}_1,..,\handwritten{\mu}_q)$ sont déterminés par récurrence à partir des relations (15), les séries entières ayant ces coefficients seront solutions du problème \uline{pourvu qu'elles soient sommables dans un\unsure{voisina}\handwritten{-}ge de l'origine} dans $K^{q+1}$ ; . En effet, s'il en est ainsi, ces séries sont fonctions continues dans un voisina ge de l'origine, donc on peut trouver un tel voisinage $U\handwritten{\times}W$ dont l'image par $(x;z_1,..,z_q)\handwritten{\to}(u_i(x;z_1,\ldots,z_q))$ soit contenue dans P, et tel que les séries $u_i$ soient substituables dans (14) pour $x\handwritten{\in}U$, $(z_1,\ldots,z_q)\handwritten{\in}W$ (\handwritten{\S}1,th.2) ; l'expression de la dérivée d'une série entière et les relations (15) montreront alors que les $u_i$ sont bien solutions du système (12) dans $U\handwritten{\in}W$.

Pour prouver la somm bilité des séries $u_i$, considérons pour chacune des séries (14), une série majorante, dont nous désignerons les coefficients par $G_i(n;r_1,..,r_p;\handwritten{\lambda}_1,\ldots,\handwritten{\lambda}_q)$ ; déterminons alors les coefficients $H_i(m;\handwritten{\mu}_1,\ldots,\handwritten{\mu}_q)$ par les conditions
